\section{Related Work}

\citet{eger2019visual} show that visually similar character substitutions substantially damage neural and non-neural NLP systems across character-, word-, and sentence-level tasks. They also study visual embeddings, adversarial training, and rule-based recovery. \citet{boucher2022bad} broaden the threat model to imperceptible encoding-level attacks, including invisible characters, homoglyphs, reorderings, and deletion controls, and demonstrate black-box failures in deployed and open-source systems.

Both studies motivate a representation-level comparison, but their scope is broader than the present controlled classification setting. Eger et al.\ provide VIPER attack and shielding code, while Boucher et al.\ provide generators and detectors for encoding perturbations. Here, those ideas are narrowed to auditable, single-code-point transformations with fixed budgets so every classifier can be evaluated on identical rows.

Tokenization provides a complementary explanation. \citet{chai2024tokenization} analyze subword sensitivity to typographical and formatting variation and show that larger models do not eliminate the problem. \citet{geh2025adversarial} manipulate the selection of valid token sequences without changing the visible input. Their attack is not a Unicode perturbation, but it establishes tokenization itself as a security-relevant interface.

Chai et al.\ release their TKEval code and data and report that regularized tokenization can improve resilience. Geh et al.\ instead exploit alternative valid segmentations at the model interface; this is adjacent tokenization-security work, not evidence about Unicode homoglyphs. Our token edit, count, and overlap measurements connect encoding perturbations to failures without conflating these distinct threat models.

For defenses, \citet{cooper2025lies} use proprietary and open-source language models to normalize homoglyph-obfuscated text. Our focus is deliberately cheaper: deterministic normalization, removal of default-ignorable characters, and a restricted confusable map. Unicode normalization has precise equivalence guarantees, but NFKC can erase meaningful compatibility distinctions \citep{unicode15}. Likewise, confusable skeletons are intended for security analysis rather than arbitrary display-text replacement \citep{unicode39}. We therefore measure clean-text harm and present confusable repair as an English-specific heuristic.

Cooper et al.\ release code and an augmented tweet dataset, enabling comparison with more flexible but more expensive normalization. Our experiment does not claim state of the art across all adversarial text defenses. Instead, it establishes transparent low-cost baselines whose transformations can be inspected code point by code point and whose clean-input behavior can be measured directly.
